\section{既有实验结果的综合}
\label{sec:prior-findings}

条件比较不仅保存数值，也能识别不同研究反复出现的规律。综合目标相和相关化合物的已发表实验，可得到三项直接服务于下一轮实验的认识。

\subsection{组成坐标比单一最高温度更能解释成相差异}

Ba--Y--Si--O助熔研究的151组实验显示，Y$_2$O$_3$/SiO$_2$比和MoO$_3$用量会主导硅酸根连接方式与晶体类型。降低Y$_2$O$_3$/SiO$_2$比或在一定范围内增加MoO$_3$，结构均趋向更高的SiO$_4$连接度；MoO$_3$过量则不利于晶体生长，较慢冷却通常改善晶体形貌和产量\cite{wierzbickawieczorek2017high}。因此，目标相筛查应把营养物组成、助熔剂比例和冷却速率作为独立变量，不能只复制一个最高温度。早期伴生晶体在1150\,$^\circ$C保温并慢冷后获得，但该结果来自含K$_2$CO$_3$/MoO$_3$的整批助熔体系，不能换算为目标相的化学计量投料\cite{kolitsch2009crystal}。

固相研究也指向同一结论。$\mathrm{Ba_xY_{26}Si_{16}O_{71+x}}$陶瓷在$x=9$--14组成系列中出现随名义Ba含量变化的主相和副相；玛瑙研磨带入的SiO$_2$会干扰本征相区判断\cite{yamane2024microstructure}。$\mathrm{Ba_5Y_{13}[SiO_4]_8O_{8.5}}$是在跨相图组成的粉体筛查中定位，并经重复退火、研磨及多种衍射和成分方法确认\cite{gulay2024navigation}。因此，目标相目前缺少的不是另一个孤立温度，而是Ba--Y--Zn--Si--O局部相图，以及与每个组成坐标对应的定量相分析。

\subsection{单晶存在不等于建立了批量相纯窗口}

目标相的直接论文用单晶衍射确认了结构，但现有资料没有同时给出可独立复现的投料、热史和批量相分数\cite{ababaikeri2024ba5y12zn}。无Zn谱系的陶瓷研究则表明，即使形成主相，批量样品仍可能含副相和处理介质污染\cite{yamane2024microstructure,gulay2024navigation}。因此，复现实验至少需要两类互补表征：单晶或先进衍射确认结构模型，定量PXRD/Rietveld测定整批样品中的目标相比例。二者缺一时，“获得目标晶体”不能等同于“建立相纯制备方法”。

\subsection{相图是连接复现与发现的核心实验}

$\mathrm{Ba_5Y_{13}[SiO_4]_8O_{8.5}}$的发现表明，即使BaO--Y$_2$O$_3$--SiO$_2$相区已有多种化合物，新相仍可能位于竞争相或玻璃区边界。研究者根据总体组成、已知相比例和成分分析逐轮更新下一批投料，才逐步缩小目标区域\cite{gulay2024navigation}。含Zn体系更需要这种策略，因为Zn的进入会同时改变阳离子数、氧计量和可能的液相行为。相图实验可区分三种情形：窄组成线化合物、无Zn谱系中的有限固溶范围，或只沿特定液相路径形成的亚稳晶相。与反复尝试单一温度相比，这三种结果都能提供更多结构信息。
